\documentclass[11pt]{article}
\usepackage[margin=1.1in]{geometry}
\usepackage{amsmath,amssymb,amsthm,mathtools}
\usepackage{booktabs,enumitem,url,xcolor}
\usepackage[colorlinks=true,linkcolor=blue!50!black,citecolor=blue!50!black,urlcolor=blue!50!black]{hyperref}
\newcommand{\E}{\mathbb{E}}
\newcommand{\PP}{\mathbb{P}}
\newcommand{\R}{\mathbb{R}}
\newcommand{\one}{\mathbf{1}}
\newcommand{\Fc}{\mathcal{F}}
\newcommand{\hm}{\mathfrak{h}}
\DeclareMathOperator{\KL}{D_{\mathrm{KL}}}
\theoremstyle{plain}
\newtheorem{theorem}{Theorem}[section]
\newtheorem{proposition}[theorem]{Proposition}
\newtheorem{lemma}[theorem]{Lemma}
\newtheorem{corollary}[theorem]{Corollary}
\theoremstyle{definition}
\newtheorem{definition}[theorem]{Definition}
\newtheorem{hypothesis}[theorem]{Hypothesis}
\theoremstyle{remark}
\newtheorem{remark}[theorem]{Remark}
\title{Conditioning a Record on Its Own Future:\\ \large Targeted Accordance and Survivorship-Conditioned Path Measures}
\author{Ajax Benander\thanks{Email:
\href{mailto:akb9408@rit.edu}{akb9408@rit.edu}. The survivorship model and
its $Q$-process limit are proved in a companion in applied probability
\cite{SurvA} and imported here; the attribute-level apparatus descends from
the Nexic companion \cite{Nexus}. The empirical hypothesis that the record of
natural history is drawn under the survivorship-conditioned law, and its
consequences, are developed in a separate downstream companion on the
Route World hypothesis, not cited here so that the dependency runs one
way, and the theorems here and the hypothesis there carry their
different grades separately. Drafting collaboration: Anthropic Claude (Opus~4.x,
Fable~5), under the author's direction.}}
\date{September 2026 --- Working Draft, rev.~4}
\begin{document}
\maketitle

\begin{abstract}
A future condition can reshape the statistical profile of earlier events
without anything in the future acting on the past: conditioning a law over
histories on a future event reweights the entire path, and the Doob
transform is the standard form of that reweighting. This paper states the
attribute-level consequence and keeps it separate from the inference the
consequence does not license. The claim in one sentence: if the realized
member of a binary partition of the past is drawn under a law conditioned
on a future target, then for every $\kappa>1$ the realized member is, with
probability at least $1-1/(1+\kappa)$, within a $\kappa$-fold exceedance
of being the more target-productive member, a statement indifferent to the
target's tense; and whether a given record is drawn under such a law is a
separate hypothesis, tested by comparing measures. Four things follow at
their stated grades. The theorem, with its one-line proof. Its
specialization to survival targets, where the critical regime forbids
conditioning on a null event and the right object is the horizon family
$\{T>t\}$ with the $Q$-process as its limit. The between-model layer: the
Bayes factor between a survivorship-conditioned and an unconditioned law,
which at finite information is the Doob density itself, so that a
historical model comparison and a stochastic $h$-transform are one object;
the realized log-evidence fluctuates and the monotone quantity is the
relative entropy, whose rate is the information flux. And the limit on
certification: the two laws are mutually absolutely continuous on every
finite horizon and singular only at infinity, so no event of bounded
duration separates them with probability one, while statistical
discrimination remains available and is what the Bayes factor measures.
A conjunction of predeclared attributes sharpens the theorem
multiplicatively under submultiplicativity, and the multiplicity guard
that any retrospective scan owes is stated as part of the result.
\end{abstract}

\section{Introduction}\label{sec:intro}

Conditioning a stochastic process on a future event is standard. The
process conditioned to avoid absorption indefinitely, the $Q$-process, is
constructed as a limit of finite-horizon conditionings and realized as a
Doob $h$-transform by the persistence-harmonic function; the penalization
literature treats a broad class of such conditionings uniformly
\cite{RVY06,MV12,CV16}. What is standard in probability has a
consequence in the interpretation of records that is not usually drawn,
and this paper draws it at the level of attributes. If a history is
sampled under a law conditioned on a future target, then each realized
attribute of that history, read against its complement, is constrained by
the target in a way the unconditioned law does not impose: the realized
member is typically the more target-productive of its pair, and a
quantitative bound says how typically. Nothing in the bound requires the
future to act on the past. The conditioning is a statement about which
histories are in the support and with what weight, and the past of a
weighted history carries the weight's fingerprints.

The paper's discipline is to keep two statements apart that are easy to
run together. The first is a theorem: \emph{under} a target-conditioned
draw, realized attributes accord with the target. The second is an
inference: \emph{because} a realized attribute accords with a target, the
record was drawn under conditioning on it. The second does not follow
from the first. It is a comparison between rival laws for the record, and
it is answered by a Bayes factor, not by the theorem. Both are stated
here, in that order, and the theorem's proof is placed beside the
disclaimer of the inference so that the two cannot be confused.

\emph{The claim, in one sentence.} If the realized member of a binary
partition of the past is drawn under a law conditioned on a future
target, then for every $\kappa>1$ it is, with probability at least
$1-1/(1+\kappa)$, within a $\kappa$-fold exceedance of being the more
target-productive member, a statement indifferent to the target's tense;
whether a given record is drawn under such a law is a separate hypothesis,
tested by comparing measures.

\emph{What is imported.} The survivorship companion \cite{SurvA} proves
the model in which survival targets are specialized here, the
finite-horizon conditioned laws and their $Q$-process limit in the
critical regime, and the theorem that the limit is locally absolutely
continuous with respect to the unconditioned law on survivors and
globally singular. The Nexic companion \cite{Nexus} supplies the
attribute-level setting: an information state as a set of stated
attributes with the unstated blank, the postulate that a realized
attribute is drawn by its parent-conditional weights, and the exceedance
modality. Each is restated in Section~\ref{sec:setting} at its grade.

\emph{What is new.} The theorem in its general form
(Section~\ref{sec:theorem}), stated for an arbitrary target and proved
in one line, with its prior-bound corollary in joint-event form and its
conjunction form under the correct dependence assumption. The
specialization to survival targets through the horizon family
(Section~\ref{sec:survival}), which removes the zero-probability
conditioning that a naive statement of the survival target commits. The
between-model layer (Section~\ref{sec:models}): the Bayes factor between
conditioned and unconditioned laws identified with the Doob density at
finite information, and the information flux as the monotone statistic
the realized evidence lacks. The limit on certification
(Section~\ref{sec:ceiling}) stated at its exact strength. And the
protocol under which the theorem's application to a real record is a
prospective test rather than a retrospective fit
(Section~\ref{sec:protocol}).

\emph{Position in the branch.} This paper's core theorem
(Sections~\ref{sec:setting}--\ref{sec:theorem}) depends on nothing
Golden and on no sibling; it is a statement about arbitrary target
conditioning. Its survival specialization (Section~\ref{sec:survival})
imports the survivorship companion \cite{SurvA}, which is the root of
the branch; a downstream technoanatomy companion imports this paper and the
survivorship one for its mechanism; and a terminal Route World
companion imports all three to state its hypothesis. Neither downstream
paper is cited here, so that the citation graph, like the logical one,
runs one way. No downstream paper is cited here as support for
anything.

\emph{What is refused.} No claim that any actual record is drawn under a
future-conditioned law is made here; that is the downstream companion's
hypothesis, and it is carried there at a hypothesis's grade. No
retrocausal mechanism is invoked anywhere. And no result below certifies
anything about a record from finite data, which Section~\ref{sec:ceiling}
proves impossible.

\section{Setting}\label{sec:setting}

\begin{definition}[Laws over histories, and partitions of the past]
\label{def:setting}
Let $(\Omega,\Fc,\PP)$ carry a law over histories, with a filtration
$(\Fc_t)$ recording the history to time $t$. An \emph{attribute} is an
$\Fc_\infty$-measurable event; a \emph{binary partition} of the past by
an attribute is a pair $\{A,A'\}$ with $A'=\Omega\setminus A$ and
$0<\PP(A)<1$. The \emph{realized member} $A^\ast$ is the member of the
pair that obtains, and $A^{\ast\prime}$ its complement. A \emph{target} is
an attribute $B$ with $\PP(B)>0$, and the \emph{$B$-conditioned law} is
$\PP_B:=\PP(\,\cdot\mid B)$. Nothing in the definition orders $A$ and $B$
in time: an attribute of the history's future is an attribute.
\end{definition}

\begin{definition}[Information states and the draw; from \cite{Nexus}]
\label{def:draw}
An information state fixes a set of stated attributes and treats the
unstated as blank rather than as unspecified. The \emph{draw} is the
postulate that at each open attribute the realized member is distributed
by the parent-conditional weights of the law in force; under $\PP$ the
weights are $\PP(A),\PP(A')$, and under $\PP_B$ they are
$\PP(A\mid B),\PP(A'\mid B)$. The modality $\lesssim$ used below is the
\emph{exceedance form}: $X\lesssim Y$ abbreviates the family of statements
that $X\ge\kappa Y$ has probability at most $1/(1+\kappa)$ for every
$\kappa>1$, the probability taken over the draw.
\end{definition}

\begin{definition}[The survivorship model; from \cite{SurvA}]
\label{def:model}
A history carries a hazard process $h_t>0$ and a killing time $T$ with
$\PP(T>t\mid\Fc^{h}_\infty)=e^{-\Lambda_t}$, $\Lambda_t=\int_0^th$. The
finite-horizon conditioned laws are $Q_t:=\PP(\,\cdot\mid T>t)$. In the
\emph{critical regime}, $\PP(T=\infty)=0$ while $\E[T]=\infty$, and the
$Q$-process is the limit $Q:=\lim_{t\to\infty}Q_t$ on each $\Fc_s$.
Everything in this paper is written for the abstract density
\[
Z_s:=\frac{dQ}{d\PP}\Big|_{\Fc_s},
\]
and two specializations are imported at their grades: the hard-wall
law, proved, $Z^{\mathrm{hard}}_s=\one_{\{\tau>s\}}\hm(X_s)/\hm(X_0)$ with
$\hm$ the persistence-harmonic function; and the soft-wall law,
conditional on the companion's ratio hypothesis (its Ansatz B),
$Z^{\mathrm{soft}}_s=e^{-\int_0^sq(X_u)du}\,g(X_s)/g(X_0)$, which
coincides with the hard-wall density exactly when $g$ is proportional to
$\hm$ after a credit translation. No result below depends on which
specialization is in force.
\end{definition}

\begin{theorem}[Locally ours, globally elsewhere; \textbf{[P]} in \cite{SurvA}]
\label{thm:singular}
In the critical regime $Q|_{\Fc_s}\ll\PP|_{\Fc_s}$ for every finite $s$,
with mutual absolute continuity on the survival event, while $Q\perp\PP$
on $\Fc_\infty$.
\end{theorem}

\section{Future-targeted accordance}\label{sec:theorem}

\begin{theorem}[Future-targeted accordance]\label{thm:accordance}
Let $\{A,A'\}$ be a binary partition of the past and $B$ a target, and
suppose the realized member $A^\ast$ is drawn under $\PP_B$. Then for every
$\kappa>1$,
\begin{equation}\label{eq:accordance}
\PP_B\!\left[\frac{\PP(B\mid A^{\ast\prime})}{\PP(B\mid A^\ast)}
\;\ge\;\kappa\,\frac{\PP(A^\ast)}{\PP(A^{\ast\prime})}\right]
\;\le\;\frac{1}{1+\kappa}.
\end{equation}
Equivalently, under the $B$-conditioned draw,
$\PP(B\mid A^{\ast\prime})/\PP(B\mid A^\ast)\lesssim\PP(A^\ast)/\PP(A^{\ast\prime})$.
\end{theorem}

\begin{proof}
By Bayes, for either member $C$ of the pair,
$\PP(B\mid C')/\PP(B\mid C)=\bigl(\PP(C'\mid B)/\PP(C\mid B)\bigr)\cdot
\bigl(\PP(C)/\PP(C')\bigr)$, so the event in \eqref{eq:accordance} is
$\{\PP(A^{\ast\prime}\mid B)/\PP(A^\ast\mid B)\ge\kappa\}$, that the
realized member carries $B$-conditioned weight $w:=\PP(A^\ast\mid B)\le
1/(1+\kappa)$. For $\kappa>1$ the threshold is below one half, so at most
one member of the pair has conditioned weight below it. Under $\PP_B$ that
member, if it exists, is realized with probability equal to its
conditioned weight, which is at most $1/(1+\kappa)$; if it does not exist
the event is empty.
\end{proof}

\begin{remark}[What the theorem says, and what it does not]
\label{rem:scope}
The outer probability in \eqref{eq:accordance} is $\PP_B$. The theorem
says that a record drawn under conditioning on $B$ has realized
attributes that typically accord with $B$: the more $B$-productive member
of each pair is typically the one realized, to within an exceedance the
draw prices. It does not say that a record whose realized attributes
accord with $B$ was drawn under $\PP_B$. That is a comparison between
$\PP_B$ and rival laws for the record, and it is Section~\ref{sec:models}.
The theorem is a within-model statement; the inference is between-model;
and the whole discipline of applying this result to an actual record is
to keep them apart. Nothing in the statement or the proof distinguishes a
target in the past from one in the future: $B$ is an attribute, and the
draw is the same draw.
\end{remark}

\begin{corollary}[The prior bound, in joint-event form]\label{cor:prior}
Fix a member $A$ of the pair with $\PP(A)=p<\tfrac12$. Under the
$B$-conditioned draw, the joint event that $A$ is realized and is the
less $B$-productive member has probability at most $p$:
\begin{equation}\label{eq:prior}
\PP_B\bigl[\,A^\ast=A\ \text{and}\ \PP(B\mid A')\ge\PP(B\mid A)\,\bigr]\;\le\;p.
\end{equation}
\end{corollary}

\begin{proof}
$\PP(B\mid A')\ge\PP(B\mid A)$ holds exactly when $w:=\PP(A\mid B)\le p$,
by the Bayes identity in the proof above with $\kappa=(1-p)/p$. Under
$\PP_B$ the member $A$ is realized with probability $w$, so the joint
event has probability $w\,\one\{w\le p\}\le p$.
\end{proof}

\noindent The corollary is the precise form of the slogan that a rare
realized attribute establishes its own significance: given the
$B$-conditioned draw, a rare member is realized-and-unproductive with
probability at most its own prior. The slogan is a within-model
diagnostic, and it is not evidence for the draw.

\begin{proposition}[Conjunction: the stockpile]\label{prop:stockpile}
Let $F_1,\dots,F_n$ be a predeclared family of attributes with priors
$p_i=\PP(F_i)$, let $F=\bigcap_iF_i$, and suppose the family is
\emph{submultiplicative}, $\PP(F)\le\prod_ip_i$, which holds with
equality under independence and holds under negative association. Then,
under the $B$-conditioned draw of the partition $\{F,F'\}$,
\begin{equation}\label{eq:stockpile}
\PP_B\bigl[\,A^\ast=F\ \text{and}\ \PP(B\mid F')\ge\PP(B\mid F)\,\bigr]
\;\le\;\PP(F)\;\le\;\prod_ip_i:
\end{equation}
the joint event that the conjunction is realized and is the less
$B$-productive member of its pair has probability at most the product
of the priors, which tightens multiplicatively in $n$ while the
complement, the union of the $F_i'$, accumulates weight $1-\prod_ip_i$.
Positive association gives $\PP(F)\ge\prod_ip_i$ and defeats the bound;
it is not a sufficient condition.
\end{proposition}

\begin{proof}
Corollary~\ref{cor:prior} applied to the fixed member $F$ of the
partition $\{F,F'\}$ gives the first inequality when $\PP(F)<\tfrac12$;
the second is submultiplicativity.
\end{proof}

\begin{proposition}[The multiplicity guard: expectation and tail]
\label{prop:guard}
If $n$ partitions are interrogated under the $B$-conditioned draw and
$N$ counts those whose realized member exhibits a $\kappa$-fold
exceedance of \eqref{eq:accordance}, then $\E[N]\le n/(1+\kappa)$; and,
with no assumption on the dependence among the partitions,
\begin{equation}\label{eq:tail}
\Pr[N\ge m]\;\le\;\frac{n}{m(1+\kappa)},\qquad
\Pr[N\ge1]\;\le\;\frac{n}{1+\kappa},
\end{equation}
by Markov's inequality and the union bound respectively. Under
independence or a stated dependence structure the tail sharpens.
Exceeding the expectation is not significance: a declared family is
scored against the tail \eqref{eq:tail}, and a family whose size $n$ is
not declared cannot be scored at all.
\end{proposition}

\begin{proof}
Linearity of expectation over the $n$ indicators, each of probability
at most $1/(1+\kappa)$ by Theorem~\ref{thm:accordance}; Markov on $N$;
and the union bound on the $n$ events.
\end{proof}

\section{Survival targets and the horizon family}\label{sec:survival}

The target of most interest is indefinite survival, and in the critical
regime it cannot be conditioned on directly: $\{T=\infty\}$ is null, and
$\PP(\,\cdot\mid T=\infty)$ is undefined. The right objects are the
horizon events.

\begin{theorem}[Horizon-indexed accordance; the second statement at the
grade of the specialization in force]\label{thm:horizon}
Let $B_t=\{T>t\}$ and $Q_t=\PP(\,\cdot\mid B_t)$. For every finite $t$ and
every binary partition $\{A,A'\}$ of $\Fc_s$, $s<t$, if the realized member
is drawn under $Q_t$ then \eqref{eq:accordance} holds with $B=B_t$ and
$\PP_B=Q_t$. In the critical regime, as $t\to\infty$ the conditioned law
$Q_t$ converges to the $Q$-process on $\Fc_s$, and for each cell
\begin{equation}\label{eq:cellconv}
\frac{\PP(T>t\mid A)}{\PP(T>t)}\;\longrightarrow\;\frac{Q(A)}{\PP(A)}
\;=\;\E_\PP[Z_s\mid A],
\end{equation}
so that
\begin{equation}\label{eq:ratioconv}
\frac{\PP(T>t\mid A')}{\PP(T>t\mid A)}\;\longrightarrow\;
\frac{Q(A')/\PP(A')}{Q(A)/\PP(A)}.
\end{equation}
\end{theorem}

\begin{proof}
The first statement is Theorem~\ref{thm:accordance} with the target
$B_t$. For \eqref{eq:cellconv},
$\PP(T>t\mid A)/\PP(T>t)=\E_\PP[\one_A\,\PP(T>t\mid\Fc_s)/\PP(T>t)]/\PP(A)$,
and by Definition~\ref{def:model} $\PP(T>t\mid\Fc_s)/\PP(T>t)\to Z_s$
in $L^1(\PP)$ on the survival event, giving $\E_\PP[Z_s\one_A]/\PP(A)
=Q(A)/\PP(A)$; \eqref{eq:ratioconv} is the quotient of two instances.
\end{proof}

\begin{remark}[Grade propagation in this section]\label{rem:grades}
Every statement in this section that consumes the survivorship
companion's conditioned law inherits that law's grade: with the
hard-wall specialization the statements are at \textbf{[P]}; with the
soft-wall specialization they are conditional on the companion's ratio
hypothesis, Ansatz~B, and carry that conditionality wherever they are
used. Theorem~\ref{thm:accordance} and its corollaries depend on neither
and remain unconditional.
\end{remark}

\begin{remark}[The ideal limit is a limit]\label{rem:limit}
Every statement about ``conditioning on indefinite survival'' in this
paper and its companions is to be read through Theorem~\ref{thm:horizon}:
a statement at finite horizon, and then a limit. The $Q$-process is not a
conditional probability on a null event, and expressions of the form
$\PP(G\mid A)$ with $G$ the event of indefinite survival are shorthand for
$\lim_t\PP(B_t\mid A)$ where that limit exists. In the survivorship model
the persistence-harmonic function $\hm$ is what the limit produces, and
the attribute-level ratios are ratios of its averages.
\end{remark}

\section{Rival laws for a record: the between-model layer}
\label{sec:models}

\begin{definition}[Hypotheses about the outer measure]\label{def:hyps}
For a record with law $\PP$ and a target family $B_t$, let $H_\PP$ be the
hypothesis that the record's realized attributes are drawn under $\PP$
(no conditioning), and $H_Q$ the hypothesis that they are drawn under
$Q$ (the survivorship-conditioned law, through its horizon family). Other
targets, an observer-existence condition among them, define other
hypotheses in the same way.
\end{definition}

\begin{proposition}[The Bayes factor is the Doob density]\label{prop:bf}
For a predeclared record statistic $E\in\Fc_s$, the Bayes factor between
the conditioned and unconditioned laws is
\begin{equation}\label{eq:bf}
\mathrm{BF}_{Q:\PP}(E)=\frac{Q(E)}{\PP(E)},
\end{equation}
and on the full information $\Fc_s$ it is the likelihood ratio $Z_s$ of
Definition~\ref{def:model}, the hard-wall Doob density or the soft
Feynman--Kac density according to the specialization in force. If $E$ is required by the
target, $Q(E)\approx1$, and rare under $\PP$, $\PP(E)=q$, then
$\mathrm{BF}_{Q:\PP}(E)\approx1/q$.
\end{proposition}

\begin{proof}
\eqref{eq:bf} is the definition of the Bayes factor for a simple event;
on $\Fc_s$ the factor is the Radon--Nikodym density by definition, and the
density's form is Definition~\ref{def:model}.
\end{proof}

\begin{remark}[The null for a surviving observer is the survivor law,
not raw $\PP$]\label{rem:survnull}
Under raw $\PP$ there is positive mass on $\{T\le s\}$, where
$Z_s=0$ and $\log Z_s=-\infty$. That is correct as mathematics and
wrong as a null for an observer whose record already exists at $s$:
current survival is evidence both hypotheses share, and a comparison
that leaves it in credits the conditioned law for the trivial fact that
the observation was made. The null for the observed surviving cohort is
$\PP^{\mathrm{surv}}_s:=\PP(\,\cdot\mid T>s)$, against which the
likelihood ratio on $\Fc_s$ is
\begin{equation}\label{eq:survlr}
\frac{dQ}{d\PP^{\mathrm{surv}}_s}\Big|_{\Fc_s}=\PP(T>s)\,Z_s,
\end{equation}
which in the hard-wall specialization is $\PP(T>s)\,\hm(X_s)/\hm(X_0)$,
finite on the whole cohort. More generally every rival law in a
comparison is conditioned on the common baseline $C$ the record fixes,
observers existing now and survival to the observation time, and the
Bayes factor is $\PP(E\mid C,H_1)/\PP(E\mid C,H_2)$, so that no part of
it is earned by evidence both rivals already had to explain.
\end{remark}

\begin{remark}[Why this and not the theorem carries the inference]
\label{rem:whybf}
A low-probability realization under $H_\PP$ is not a contradiction of
$H_\PP$; it is evidence against it, and \eqref{eq:bf} says how much. The
attribute-level accordance of Theorem~\ref{thm:accordance} tells one what
to expect of a record if $H_Q$ holds; the Bayes factor tells one whether
$H_Q$ holds better than a rival. A record can accord with $Q$ at every
interrogated attribute and still be better explained by a rival law that
also predicts those attributes, which is why the comparison must name its
rival and predeclare its statistic. The same object, the Doob density,
serves both layers: within the model it is the weight by which the
conditioning tilts each path, and between models it is the evidence the
tilted path carries.
\end{remark}

\begin{definition}[Realized evidence, information, and flux]\label{def:info}
Let $Z_t=dQ|_{\Fc_t}/d\PP|_{\Fc_t}$ and $\ell_t=\log Z_t$, the realized
log-evidence. Since $Z$ is a $\PP$-martingale, $\ell_t$ is not monotone
along a path: new observations can favor $Q$, then $\PP$. The monotone
quantity is the \emph{information}
\begin{equation}\label{eq:info}
I_t:=\E_Q[\ell_t]=\KL\bigl(Q|_{\Fc_t}\,\big\|\,\PP|_{\Fc_t}\bigr),
\end{equation}
nondecreasing in $t$ as the filtration grows, with
$\E_\PP[\ell_t]=-\KL(\PP|_{\Fc_t}\|Q|_{\Fc_t})\le0$. Where differentiable,
$\mathcal J_t:=dI_t/dt$ is the \emph{information flux}.
\end{definition}

\begin{proposition}[Discriminating rival laws by the law of the evidence]
\label{prop:discriminate}
Under $H_Q$ the realized log-evidence has mean $I_t\ge0$ and, in the
critical regime of the survivorship model, $I_t\to\infty$; under $H_\PP$
it has mean $-\KL(\PP_t\|Q_t)\le0$. The two hypotheses are therefore
discriminated by the distribution of $\ell_t$, of which the flux is the
summary, and not by any monotone accumulation along a single path, which
does not exist.
\end{proposition}

\begin{proof}
The means are \eqref{eq:info} and its $\PP$-counterpart; divergence under
$H_Q$ in the critical regime follows from Theorem~\ref{thm:singular},
since mutual singularity on $\Fc_\infty$ forces the relative entropy to
diverge.
\end{proof}

\section{The limit on certification}\label{sec:ceiling}

\begin{theorem}[No bounded observation certifies the law]\label{thm:ceiling}
In the critical regime, for every finite $s$, no event in $\Fc_s$ has
probability one under $Q$ and probability zero under $\PP$ on the
survival event, nor conversely. Consequently no observation of bounded
duration certifies whether a surviving record is drawn under $Q$ or under
$\PP$; certification is available only on $\Fc_\infty$, where the laws are
singular.
\end{theorem}

\begin{proof}
Mutual absolute continuity on the survival event at every finite $s$ is
Theorem~\ref{thm:singular}; an event of probability one under one law and
zero under the other would violate it.
\end{proof}

\begin{remark}[Certification against discrimination]\label{rem:certify}
The theorem forbids perfect separation and nothing more. Mutually
absolutely continuous laws can be statistically far apart, and
Proposition~\ref{prop:bf} is exactly the measure of how far: a Bayes
factor of $10^{8}$ on a finite window is discrimination of a high order,
and it is not certification, because the same window has positive
probability under both laws. The theorem is turned on this paper's own
subject matter deliberately: whatever a record's attributes say in favor
of $H_Q$, they say it at a Bayes factor and never at probability one.
\end{remark}

\section{A protocol for applying the theorem to a record}
\label{sec:protocol}

The theorem applies to any record with a law and a target. Its
application to an actual record faces the retrospective-target objection
in its full force: any realized history, read backward, contains rare
attributes that are prerequisites of whatever it happens to contain, so a
set of attributes selected after the fact and found to accord with a
target chosen after the fact establishes nothing. Two of the objection's
three freedoms are closed by the theorem's form: the target is fixed
before the record is read, and the alternative to each attribute is its
complement, not an analyst's choice of class. The third freedom, which
attributes to interrogate, is not closed by any theorem, and the
multiplicity guard prices it only when the scan is declared. The protocol
is therefore prospective.

\begin{definition}[Prospective test of a target hypothesis]\label{def:protocol}
Declare, before the observation interval: (i) the target family $B_t$ and
the rival law; (ii) the attribute family $F_1,\dots,F_n$, with $n$ fixed;
(iii) the baseline priors $p_i$ under the rival law and the dependence
assumption on the family (independence, or negative association, or a
stated bound on $\PP(\bigcap F_i)$); (iv) the record statistic $E$ whose
Bayes factor \eqref{eq:bf} will be reported; (v) the stopping rule and
the multiplicity budget $n/(1+\kappa)$ at the $\kappa$ to be scored. Then
observe. A retrospective exercise on an existing record, however
suggestive, is hypothesis-generating and is reported as such.
\end{definition}

\section{Open problems}\label{sec:open}

Every open item is gathered here. The soft-wall ratio hypothesis on which
the survivorship model's $Q$-process rests is open in the companion
\cite{SurvA}, and Theorem~\ref{thm:horizon}'s second statement inherits
it. The convergence of $Q_t$ on a window shifted to the present, the law
of $(X_{t+u})_{u\le s}$ given survival to $t$, is a different statement
from convergence on an early window and is not established; every claim
about a surviving record's \emph{future} behavior waits on it. A finite,
monotone evidence budget under which a Copernican position statement
would be well posed is not defined, and no Lindy-type result is claimed.
And the extension of the multiplicity guard from a fixed family to a
sequentially chosen one, with a stopping rule, is stated as a protocol
requirement rather than proved as a bound.

\section{Conclusion}\label{sec:conclusion}

The claim stated at the outset stands as proved. If the realized member
of a binary partition of the past is drawn under a law conditioned on a
future target, then for every $\kappa>1$ it is, with probability at least
$1-1/(1+\kappa)$, within a $\kappa$-fold exceedance of being the more
target-productive member, and the statement is indifferent to the
target's tense; whether a given record is drawn under such a law is a
separate hypothesis, tested by comparing measures, and the measure of
that comparison is the Doob density. The future does nothing to the past
in any of this. A conditioned law has a support and a weight, and the
past of a weighted history carries the weight; that is the whole of what
``conditioning a record on its own future'' means, and it is enough to
make the companion's hypothesis about the record of natural history a
hypothesis with a stated test rather than a reading with a stated
feeling.

\begin{thebibliography}{99}\small
\bibitem{SurvA} A.~Benander, \emph{Survivorship conditioning for hazard
rates driven at the log-derivative level}, working draft, 2026.
\bibitem{Nexus} A.~Benander, \emph{Nexic Reasoning}, companion manuscript,
2026.
\bibitem{RVY06} B.~Roynette, P.~Vallois, and M.~Yor, \emph{Some penalisations
of the Wiener measure}, Japanese J.\ Math.\ \textbf{1} (2006), 263--290.
\bibitem{MV12} S.~M\'el\'eard and D.~Villemonais, \emph{Quasi-stationary
distributions and population processes}, Probab.\ Surveys \textbf{9} (2012),
340--410.
\bibitem{CV16} N.~Champagnat and D.~Villemonais, \emph{Exponential
convergence to quasi-stationary distribution and $Q$-process}, Probab.\
Theory Relat.\ Fields \textbf{164} (2016), 243--283 (AUTHOR TO VERIFY).
\end{thebibliography}
\end{document}
